\begin{center}
    \LARGE \textbf{CIUDAD VICIO}  \\
    \vspace{1em}
    \small Primer Escena: Estafa cambio de pañuelo
\end{center}
\raggedright

\vspace{2em}

\textbf{EXT. CALLEJÓN - DÍA}

\begin{scenedesc}
Miguel sale de una puerta pequeña al fondo de un callejón  y gira a la izquierda adentrándose en una vía más amplia. Aparentemente, la calle está vacía, cuando se escuchan unas ruedas y aparece un joven sobre una patineta en dirección contraria a la de Miguel. El joven se detiene justo a su lado, y ambos se miran durante un instante. 

De repente, Miguel oye gritos procedentes de algún lugar a sus espaldas. Un hombre de mediana edad y complexión menuda dobla la esquina a toda velocidad, corriendo desesperadamente. Lleva una cartera en la mano. A lo lejos se escuchan unos gritos.
\end{scenedesc}

\begin{dialogue}{Desconocido (Teódulo)}
¡Cabrón! No te quieras pasar de verga. Esa es mi Mochila. 
\end{dialogue}

\begin{scenedesc}
Miguel se queda inmóvil, sin el menor interés en ejercer la justicia. Justo cuando el ladrón está a punto de pasar de largo, Amaury alza su tabla y se la lanza a las piernas, haciéndolo tropezar y rodar por el suelo, provocando que suelte la Mochila que llevaba.

Amaury corre y patea la mochila en la dirección de Miguel, casi por reflejo, Miguel la recoge. El ladrón se pone en pie y saca una navaja, Miguel y Amaury se alertan en posición defensiva.

El ladrón, al darse cuenta que buscan la pelea lanza un navajaso hacia Amaury pero Amaury lo esquiva con facilidad. Después de esto, el ladrón decide huir.  
\end{scenedesc}

\begin{scenedesc}
El ladrón huye hacia la calle principal y desaparece de la vista. De entre unos vochecitos viejos y abandonados entre cristales rotos y metales puntiagudos, un viejo se arrastra hacia el medio de la calle. 

Miguel al verlo que está sangrando de la pierna lo toma de las manos y le ayuda a sentarse en el suelo mientras sostiene la mochila en el hombro. Amaury se acerca al viejo con indiferencia a Miguel.  
\end{scenedesc}
\begin{dialogue}{Miguel}
\parenthetical{Sosteniendo al viejo}
¿Está bien? Déjeme darle una ayudadita Don.
\end{dialogue}

\begin{dialogue}{Hombre Viejo (Teódulo)}
 Mi mochila y mi dinero.
\end{dialogue}

\begin{dialogue}{Miguel}
\parenthetical{Entregándosela}
Tranquilo, Don. Aquí la tengo
\end{dialogue}

\begin{dialogue}{Hombre Viejo (Teódulo)}
\parenthetical{Aliviado, tomándola}
Ese pendejo me empujo contra este carro abandonado siento como que me rompí la pierna.
\end{dialogue}

\begin{dialogue}{Desconocido (Amaury)}
No se mueva jefe, mejor le hablamos a la ambulancia o a la patrulla.
\end{dialogue}

\begin{dialogue}{Hombre Viejo (Teódulo)}
¡No, la tira  no!
\end{dialogue}

\begin{scenedesc}
Teódulo abre la mochila, revelando un fajo de billetes 500 pesos. Miguel se emociona al ver el dinero fácil. Teódulo se intenta poner de pie.
\end{scenedesc}

\begin{dialogue}{Hombre Viejo (Teódulo)}
    Me tengo que ir ya.
\end{dialogue}

\begin{dialogue}{Miguel}
    ¿Cómo? No se puede ir asi jefe.  
\end{dialogue}


\begin{dialogue}{Hombre Viejo (Teódulo)}

    No lo entiendes, es ahuevo que tengo que llegar. Yo soy el encargado del puesto de aqui la esquina, esa feria es de la unión y es la segunda vez que no llevo el corte de la semana a tiempo. Hoy me mandaron un mensajito con un tipo en moto. Si no entrego este dinero en el punto antes de las 6 van a decir que me pelé con todo y feria. Me van a quebrar.

\end{dialogue}


\begin{dialogue}{Desconocido (Amaury)}
¡Viejo!, no puede ni caminar. Ni tirándole paro llegamos a tiempo.
\end{dialogue}

\begin{dialogue}{Hombre Viejo (Teódulo)}
Tengo 3 mil si me entregan el dinero. Se van a michas.
\end{dialogue}

\begin{dialogue}{Desconocido (Amaury)}
\parenthetical{vacila, visiblemente nervioso}
¡Saco! Yo creo ya me campanearon los de la cuadra. ¿Que tal que me estan esperando?  y traigo la bronca.  
\end{dialogue}

\begin{dialogue}{Hombre Viejo (Teódulo)}
Nadie va a saber que llevas.
\end{dialogue}

\begin{dialogue}{Desconocido (Amaury)}
\parenthetical{sacudiendo la cabeza}
Si quiere le ayudo a acomodarse, pero le saco al encarguito. 
\end{dialogue}

\begin{dialogue}{Hombre Viejo (Teódulo)}
\parenthetical{desesperado, girándose hacia Miguel}
¿Que me dices tu cabrón? Te llevas los 3 completos pero hazme el favor.
\end{dialogue}

\begin{dialogue}{Desconocido (Amaury)}
No manche Don, ¿por que va a confiar en él? Esta caliente esa bronca y es mucho varo. 
\end{dialogue}

\begin{dialogue}{Miguel}
Chamaco miado, a ti que te valga verga. 
\parenthetical{se dirije a Teódulo}
¿Cual punto es?
\end{dialogue}

\begin{dialogue}{Hombre Viejo (Teódulo)}
Es junto a la capilla del centro, la casa azul, tocas en el zaguán negro. Hay veinte mil ahí, y aquí tres mil para usted.
\end{dialogue}

\begin{dialogue}{Miguel}
\parenthetical{tomando el fajo de billetes y los billetes sueltos}
Pues ya esta Don, yo me encargo, usted tranquilo. 
\end{dialogue}

\begin{scenedesc}
Miguel se guarda el dinero con total naturalidad en su chamarra junto a otro sobre con dinero proveniente de la casa de cambio. 
Amaury, que había sacado un pañuelo de su bolsillo para empezar a ``vendar'' la pierna de Teódulo, se detiene y lo observa fijamente.
\end{scenedesc}

\begin{dialogue}{Desconocido (Amaury)}
En cuanto te campaneen en la esquina va a salir el clavo. 
\end{dialogue}

\begin{dialogue}{Hombre Viejo (Teódulo)}
\parenthetical{asustado de nuevo repentinamente}
¿Como crees que lo vas a entusar asi?
\end{dialogue}

\begin{dialogue}{Desconocido (Amaury)}
Aunque sea en su calcetín maestro, ¿No tiene nada para disimular?
\end{dialogue}

\begin{dialogue}{Miguel}
No
\end{dialogue}

\begin{dialogue}{Desconocido (Amaury)}
¿Algo como un trapo, un pañuelo?
\end{dialogue}

\begin{dialogue}{Hombre Viejo (Teódulo)}
A ver aqui
\parenthetical{Teódulo saca un pañuelo blanco arrugado}
\end{dialogue}

\begin{dialogue}{Desconocido (Amaury)}
Póngame aqui el dinero
\end{dialogue}

\begin{scenedesc}
Miguel vacila un momento, pero luego de unos segundos le da el fajo de billetes a Amaury. Amaury los coloca en el centro del pañuelo.
\end{scenedesc}

\begin{dialogue}{Desconocido (Amaury)}
    Ponle todo lo que traigas man, el clavo es el clavo.
\end{dialogue}

\begin{dialogue}{Hombre Viejo (Teódulo)}
\parenthetical{suplicando}
Mi pierna, trae un doctor, tu métele prisa.
\end{dialogue}

\begin{scenedesc}
Miguel, arrastrado por la urgencia del momento, saca el dinero que traía de la casa de cambio y lo deposita también dentro del pañuelo. Amaury junta rápidamente las esquinas del pañuelo y las ata con un nudo muy firme.
\end{scenedesc}

\begin{dialogue}{Desconocido (Amaury)}
\parenthetical{haciendo una demostración al deslizar el fajo dentro de sus pantalónes}
El clavo bien clavado. Ese entuse no sale y no se ve.
\end{dialogue}

\begin{scenedesc}
Amaury saca el fajo de sus pantalones , cambiando el pañuelo falso y el real,  y le entrega el paquete falso idéntico a Miguel.
\end{scenedesc}

\begin{dialogue}{Miguel}
\parenthetical{metiéndose el fajo falso en los pantalones}
Ahuevo, si se la sabe chamaco mión.
\end{dialogue}

\begin{dialogue}{Desconocido (Amaury)}
Es sin preguntar. 
\end{dialogue}

\begin{scenedesc}
Miguel echa un vistazo a las salidas del callejón, da la vuelta y se aleja caminando lo más rápido posible sin llegar a correr. \linebreak

Amaury lo observa marcharse hasta que Miguel da vuelta en la esquina y se sube a un Taxi. En el preciso instante en que lo pierde de vista, Teódulo se incorpora por completo, olvidándose instantáneamente de su falsa cojera. Amaury saca el pañuelo real que escondía tras su espalda, lo desata y sonríe ante el gigantesco botín. Ambos comparten una breve y silenciosa carcajada antes de marcharse por caminos separados. \linebreak

Miguel, por otra parte, cuando el taxi se ha alejado del lugar, confiado de su robo saca el pañuelo de su pantalón y descubre billetes hechos de recortes de papel periódico. \\
\end{scenedesc}

